\begin{bt}%[9D6H18-2]
% nguon: Tuyển tập đề tham khảo môn Toán Tuyển sinh 10 năm học 2025-2026 TP.HCM – Đề tham khảo Huyện Củ Chi (Đề 1) – Bài 1
Cho parabol $(P): y=\dfrac{x^2}{2}$.
\begin{enumerate}[a)]
	\item Vẽ đồ thị $(P)$ trên hệ trục toạ độ $Oxy$.
	\item Tìm toạ độ những điểm thuộc $(P)$ có tung độ bằng $4$.
\end{enumerate}
\loigiai{
\begin{enumerate}[a)]
	\item Bảng giá trị tương ứng của $x$ và $y$:
	\begin{center}
	\begin{tabular}{|c|c|c|c|c|c|}
	\hline
	$x$ & $-4$ & $-2$ & $0$ & $2$ & $4$ \\ \hline
	$y=\dfrac{x^2}{2}$ & $8$ & $2$ & $0$ & $2$ & $8$ \\ \hline
	\end{tabular}
	\end{center}
	Đồ thị $(P)$ là một đường parabol đi qua gốc toạ độ $O(0;0)$, nhận trục tung $Oy$ làm trục đối xứng và quay bề lõm lên phía trên.
	\item Gọi điểm thuộc $(P)$ có tung độ bằng $4$ là $M(x_0; 4)$.\\
	Vì $M \in (P)$ nên:
	$$ \dfrac{x_0^2}{2} = 4 \iff x_0^2 = 8 \iff x_0 = \pm \sqrt{8} = \pm 2\sqrt{2}. $$
	Với $x_0 = 2\sqrt{2} \implies M_1(2\sqrt{2}; 4)$.\\
	Với $x_0 = -2\sqrt{2} \implies M_2(-2\sqrt{2}; 4)$.\\
	Vậy có hai điểm thoả mãn là $M_1(2\sqrt{2}; 4)$ và $M_2(-2\sqrt{2}; 4)$.
\end{enumerate}
}
\end{bt}

\begin{bt}%[9D6H20-2]
% nguon: Tuyển tập đề tham khảo môn Toán Tuyển sinh 10 năm học 2025-2026 TP.HCM – Đề tham khảo Huyện Củ Chi (Đề 1) – Bài 2
Cho phương trình bậc hai: $3x^2+2x-3=0$.
\begin{enumerate}[a)]
	\item Chứng minh phương trình luôn có hai nghiệm phân biệt $x_1, x_2$.
	\item Không giải phương trình, hãy tính giá trị của biểu thức:
	$$ M = (x_1-2x_2)(x_2-x_1) + x_2^2. $$
\end{enumerate}
\loigiai{
\begin{enumerate}[a)]
	\item Phương trình $3x^2+2x-3=0$ có các hệ số $a=3$, $b=2$, $c=-3$.\\
	Vì $a \cdot c = 3 \cdot (-3) = -9 < 0$ nên phương trình luôn có hai nghiệm phân biệt $x_1, x_2$.
	\item Theo hệ thức Vi-ét, ta có:
	$$ \begin{cases} S = x_1+x_2 = -\dfrac{2}{3} \\ P = x_1 x_2 = \dfrac{-3}{3} = -1. \end{cases} $$
	Khai triển và thu gọn biểu thức $M$:
	$$ M = -(x_1-2x_2)(x_1-x_2) + x_2^2 = -(x_1^2 - 3x_1 x_2 + 2x_2^2) + x_2^2 = -x_1^2 + 3x_1 x_2 - x_2^2. $$
	Biến đổi theo tổng và tích các nghiệm:
	$$ M = -(x_1^2 + x_2^2) + 3x_1 x_2 = -[(x_1+x_2)^2 - 2x_1 x_2] + 3x_1 x_2 = -(x_1+x_2)^2 + 5x_1 x_2 = -S^2 + 5P. $$
	Thay $S = -\dfrac{2}{3}$ và $P = -1$ vào biểu thức $M$:
	$$ M = -\left(-\dfrac{2}{3}\right)^2 + 5(-1) = -\dfrac{4}{9} - 5 = -\dfrac{49}{9}. $$
\end{enumerate}
}
\end{bt}

\begin{bt}%[9D8H26-3]
% nguon: Tuyển tập đề tham khảo môn Toán Tuyển sinh 10 năm học 2025-2026 TP.HCM – Đề tham khảo Huyện Củ Chi (Đề 1) – Bài 3
Bảng số liệu thống kê nhiệt độ trung bình tháng và lượng mưa các tháng trong năm 2022 tại Hà Nội cho thấy: nhiệt độ cao nhất rơi vào tháng 6 ($31^\circ\text{C}$) và lượng mưa lớn nhất rơi vào tháng 8 ($340\text{ mm}$). Trong 12 tháng của năm, có 2 tháng có nhiệt độ trung bình không quá $20^\circ\text{C}$ (tháng 1 và tháng 2) và có 4 tháng có lượng mưa vượt quá $150\text{ mm}$ (tháng 5, 6, 7, 8).
\begin{enumerate}[a)]
	\item Trong năm 2022, tháng nào ở Hà Nội có nhiệt độ cao nhất và tháng nào có lượng mưa nhiều nhất?
	\item Chọn ngẫu nhiên một tháng trong năm 2022, tính xác suất của các biến cố sau:
	\begin{itemize}
		\item Biến cố $A$: ``Tháng được chọn có nhiệt độ trung bình không quá $20^\circ\text{C}$''.
		\item Biến cố $B$: ``Tháng được chọn có lượng mưa vượt quá $150\text{ mm}$''.
	\end{itemize}
\end{enumerate}
\loigiai{
\begin{enumerate}[a)]
	\item Dựa vào số liệu thống kê:
	\begin{itemize}
		\item Tháng có nhiệt độ cao nhất trong năm là tháng 6 (đạt $31^\circ\text{C}$).
		\item Tháng có lượng mưa nhiều nhất trong năm là tháng 8 (đạt $340\text{ mm}$).
	\end{itemize}
	\item Một năm có 12 tháng nên số phần tử của không gian mẫu là $n(\Omega) = 12$.
	\begin{itemize}
		\item Có $2$ tháng có nhiệt độ không quá $20^\circ\text{C}$, do đó số kết quả thuận lợi cho biến cố $A$ là $n(A) = 2$.\\
		Xác suất của biến cố $A$ là:
		$$ P(A) = \dfrac{2}{12} = \dfrac{1}{6}. $$
		\item Có $4$ tháng có lượng mưa vượt quá $150\text{ mm}$, do đó số kết quả thuận lợi cho biến cố $B$ là $n(B) = 4$.\\
		Xác suất của biến cố $B$ là:
		$$ P(B) = \dfrac{4}{12} = \dfrac{1}{3}. $$
	\end{itemize}
\end{enumerate}
}
\end{bt}

\begin{bt}%[9D6V21-3]
% nguon: Tuyển tập đề tham khảo môn Toán Tuyển sinh 10 năm học 2025-2026 TP.HCM – Đề tham khảo Huyện Củ Chi (Đề 1) – Bài 4
Một khu vườn hình chữ nhật có chiều dài $100\text{ m}$ và chiều rộng $50\text{ m}$. Người ta bớt ra một phần đất dọc theo chu vi và góc vườn để làm lối đi và khu vực trồng hoa có bề rộng phụ thuộc vào tham số $x\text{ (m)}$ ($0 < x < 25$). Sau khi hoàn thành, phần đất còn lại dành riêng cho khu vực trồng cây ăn trái có kích thước chiều dài là $(100 - 4x)\text{ m}$ và chiều rộng là $(50 - 2x)\text{ m}$.
\begin{enumerate}[a)]
	\item Viết biểu thức $S$ biểu diễn theo $x$ diện tích phần đất của khu vực trồng cây ăn trái.
	\item Tìm giá trị của $x$ biết phần đất trồng cây ăn trái có diện tích là $4608\text{ m}^2$.
\end{enumerate}
\loigiai{
\begin{enumerate}[a)]
	\item Biểu thức diện tích phần đất của khu vực trồng cây ăn trái theo $x$ là:
	$$ S(x) = (100 - 4x)(50 - 2x) = 5000 - 200x - 200x + 8x^2 = 8x^2 - 400x + 5000\text{ (m}^2\text{)}. $$
	\item Theo đề bài, diện tích khu vực trồng cây ăn trái là $4608\text{ m}^2$, ta có phương trình:
	$$ 8x^2 - 400x + 5000 = 4608 \iff 8x^2 - 400x + 392 = 0 \iff x^2 - 50x + 49 = 0. $$
	Giải phương trình bậc hai:
	Vì $a + b + c = 1 + (-50) + 49 = 0$ nên phương trình có hai nghiệm:
	$$ x_1 = 1\text{ (nhận do } 0 < x < 25\text{)}, \quad x_2 = 49\text{ (loại vì } x < 25\text{)}. $$
	Vậy giá trị cần tìm là $x = 1\text{ m}$.
\end{enumerate}
}
\end{bt}

\begin{bt}%[9H0V31-3]
% nguon: Tuyển tập đề tham khảo môn Toán Tuyển sinh 10 năm học 2025-2026 TP.HCM – Đề tham khảo Huyện Củ Chi (Đề 1) – Bài 5
Một cây bút chì có dạng hình trụ dài $180\text{ mm}$ và đường kính thân bút là $7{,}2\text{ mm}$. Phần ruột bút làm bằng than chì có dạng hình trụ có chiều dài bằng với chiều dài của thân bút và đường kính ngòi bằng $3{,}4\text{ mm}$.
\begin{enumerate}[a)]
	\item Hãy tính thể tích than chì cần dùng để làm phần ruột cho một cây bút chì khi chưa gọt (lấy $\pi \approx 3{,}14$, kết quả làm tròn đến hàng đơn vị của $\text{mm}^3$).
	\item Để sản xuất phần vỏ gỗ của bút chì, người ta dùng những thanh gỗ hình hộp có đáy là hình vuông cạnh $8\text{ mm}$ và chiều dài $185\text{ mm}$. Hỏi với $10\text{ m}^3$ gỗ chuyên dụng làm vỏ bút chì thì có thể tạo ra được bao nhiêu cây bút chì, biết rằng trong quá trình xẻ nhỏ gỗ thì lượng hao hụt chiếm $12\%$ do mùn cưa và phế phẩm?
\end{enumerate}
\loigiai{
\begin{enumerate}[a)]
	\item Bán kính của phần ruột than chì hình trụ là:
	$$ r = \dfrac{3{,}4}{2} = 1{,}7\text{ mm}. $$
	Chiều dài của ruột than chì là $h = 180\text{ mm}$.\\
	Thể tích than chì cần dùng để làm phần ruột của một cây bút chì là:
	$$ V_{\text{ruột}} = \pi r^2 h = 3{,}14 \cdot (1{,}7)^2 \cdot 180 \approx 1634\text{ mm}^3. $$
	\item Thể tích một thanh gỗ hình hộp để tạo nên một cây bút chì là:
	$$ V_1 = 8 \cdot 8 \cdot 185 = 11\,840\text{ mm}^3 = 11\,840 \cdot 10^{-9}\text{ m}^3 = 1{,}184 \cdot 10^{-5}\text{ m}^3. $$
	Vì hao hụt trong quá trình cưa xẻ gỗ chiếm $12\%$ nên tỷ lệ thể tích gỗ hữu dụng thực tế đạt được là:
	$$ 100\% - 12\% = 88\% = 0{,}88. $$
	Thể tích gỗ thực tế được dùng để tiện vỏ bút chì từ $10\text{ m}^3$ gỗ ban đầu là:
	$$ V_{\text{hữu dụng}} = 10 \cdot 0{,}88 = 8{,}8\text{ m}^3. $$
	Số cây bút chì tối đa có thể sản xuất được là:
	$$ N = \dfrac{V_{\text{hữu dụng}}}{V_1} = \dfrac{8{,}8}{1{,}184 \cdot 10^{-5}} \approx 743\,243\text{ (cây)}. $$
	Vậy có thể tạo ra được khoảng $743\,243$ cây bút chì.
\end{enumerate}
}
\end{bt}

\begin{bt}%[9D1H3-4]
% nguon: Tuyển tập đề tham khảo môn Toán Tuyển sinh 10 năm học 2025-2026 TP.HCM – Đề tham khảo Huyện Củ Chi (Đề 1) – Bài 6
Để đảm bảo dinh dưỡng trong khẩu phần ăn hàng ngày, một gia đình gồm $4$ thành viên cần ít nhất $900$ đơn vị prôtêin và $400$ đơn vị lipít mỗi ngày. Mỗi ki-lô-gam thịt bò chứa $800$ đơn vị prôtêin và $200$ đơn vị lipít, mỗi ki-lô-gam thịt heo chứa $600$ đơn vị prôtêin và $400$ đơn vị lipít. Giá của thịt bò là $240\,000$ đồng/kg và giá của thịt heo là $160\,000$ đồng/kg. Hỏi cần mua bao nhiêu tiền thịt bò và bao nhiêu tiền thịt heo để cung cấp vừa đủ $900$ đơn vị prôtêin và $400$ đơn vị lipít trong một ngày cho gia đình trên?
\loigiai{
Gọi $x$ và $y$ lần lượt là số ki-lô-gam thịt bò và thịt heo cần mua trong một ngày ($x, y > 0$, đơn vị: kg).\\
Để cung cấp vừa đủ $900$ đơn vị prôtêin và $400$ đơn vị lipít, ta có hệ phương trình:
$$ \begin{cases} 800x + 600y = 900 \\ 200x + 400y = 400 \end{cases} \iff \begin{cases} 4x + 3y = 4{,}5 \\ x + 2y = 2 \end{cases} \iff \begin{cases} x = 2 - 2y \\ 4(2 - 2y) + 3y = 4{,}5 \end{cases} \iff \begin{cases} 5y = 3{,}5 \\ x = 2 - 2y \end{cases} \iff \begin{cases} x = 0{,}6\text{ (thoả mãn)} \\ y = 0{,}7\text{ (thoả mãn)}. \end{cases} $$
Số tiền mua thịt bò là:
$$ 0{,}6 \cdot 240\,000 = 144\,000\text{ (đồng)}. $$
Số tiền mua thịt heo là:
$$ 0{,}7 \cdot 160\,000 = 112\,000\text{ (đồng)}. $$
Tổng số tiền cần mua thịt là:
$$ 144\,000 + 112\,000 = 256\,000\text{ (đồng)}. $$
}
\end{bt}

\begin{bt}%[9H9C29-3]
% nguon: Tuyển tập đề tham khảo môn Toán Tuyển sinh 10 năm học 2025-2026 TP.HCM – Đề tham khảo Huyện Củ Chi (Đề 1) – Bài 7
Cho nửa đường tròn tâm $O$ đường kính $AB$. Trên nửa đường tròn $(O)$ lấy điểm $C$ bất kỳ ($C \ne A, C \ne B$). Trên cung nhỏ $AC$ lấy điểm $M$ sao cho $MA = MC$. Hai đường thẳng $BC$ và $AM$ cắt nhau tại điểm $E$, hai đường thẳng $BM$ và $AC$ cắt nhau tại điểm $H$.
\begin{enumerate}[a)]
	\item Chứng minh $BM$ là tia phân giác của góc $\widehat{ABC}$.
	\item Chứng minh $ME^2 = MH \cdot MB$.
	\item Đường tròn ngoại tiếp tam giác $BEH$ cắt nửa đường tròn $(O)$ tại điểm $F$ ($F \ne B$). Tia $EF$ cắt đoạn thẳng $AB$ tại điểm $P$, hai đường thẳng $BM$ và $AF$ cắt nhau tại $Q$. Chứng minh $PQ \perp AB$.
\end{enumerate}
\loigiai{
\begin{enumerate}[a)]
	\item Vì $M$ thuộc nửa đường tròn $(O)$ và dây $MA = MC$ nên hai cung chắn bởi hai dây bằng nhau: cung $MA$ bằng cung $MC$.\\
	Hai góc nội tiếp $\widehat{ABM}$ và $\widehat{MBC}$ lần lượt chắn hai cung bằng nhau cung $MA$ và cung $MC$, do đó $\widehat{ABM} = \widehat{MBC}$.\\
	Vậy $BM$ là tia phân giác của góc $\widehat{ABC}$.
	\item Vì $C$ và $M$ thuộc đường tròn đường kính $AB$ nên $\widehat{ACB} = 90^\circ \implies AC \perp BE$ tại $C$, và $\widehat{AMB} = 90^\circ \implies BM \perp AE$ tại $M$.\\
	Xét tam giác $ABE$, hai đường cao $AC$ và $BM$ cắt nhau tại trực tâm $H$, do đó $EH \perp AB$.\\
	Trong tam giác vuông $AME$ vuông tại $M$, ta có $\widehat{MEH} = \widehat{MAB}$ (cùng phụ với $\widehat{ABM}$).\\
	Lại có $\widehat{MBE} = \widehat{MAB}$ (cùng chắn hai cung bằng nhau hoặc do tam giác $ABE$ cân tại $B$ vì $BM$ vừa là phân giác vừa là đường cao).\\
	Suy ra $\widehat{MEH} = \widehat{MBE}$, từ đó hai tam giác vuông $\Delta MEH$ và $\Delta MBE$ đồng dạng (g-g):
	$$ \dfrac{ME}{MB} = \dfrac{MH}{ME} \implies ME^2 = MH \cdot MB. $$
	\item Vì đường tròn ngoại tiếp tam giác $BEH$ cắt nửa đường tròn $(O)$ tại $F$ nên tứ giác $BEFH$ nội tiếp đường tròn đường kính $EH$.\\
	Suy ra $\widehat{EFB} = 90^\circ \implies EF \perp BF$.\\
	Mặt khác, $F$ thuộc nửa đường tròn $(O)$ đường kính $AB$ nên $\widehat{AFB} = 90^\circ \implies AF \perp BF$.\\
	Từ đó ba điểm $A, F, E$ thẳng hàng (hoặc $AF$ trùng $AE$), suy ra $Q$ là trực tâm của tam giác $PAB$.\\
	Do đó đường cao thứ ba kẻ từ $Q$ hạ xuống cạnh đối diện $AB$ phải vuông góc với $AB$, tức là $PQ \perp AB$.
\end{enumerate}
}
\end{bt}

\begin{bt}%[9D6H18-2]
% nguon: Tuyển tập đề tham khảo môn Toán Tuyển sinh 10 năm học 2025-2026 TP.HCM – Đề tham khảo Huyện Củ Chi (Đề 2) – Bài 1
Cho parabol $(P): y=x^2$.
\begin{enumerate}[a)]
	\item Vẽ đồ thị $(P)$ trên hệ trục toạ độ $Oxy$.
	\item Tìm toạ độ những điểm thuộc $(P)$ có tung độ và hoành độ bằng nhau.
\end{enumerate}
\loigiai{
\begin{enumerate}[a)]
	\item Bảng giá trị tương ứng của $x$ và $y$:
	\begin{center}
	\begin{tabular}{|c|c|c|c|c|c|}
	\hline
	$x$ & $-2$ & $-1$ & $0$ & $1$ & $2$ \\ \hline
	$y=x^2$ & $4$ & $1$ & $0$ & $1$ & $4$ \\ \hline
	\end{tabular}
	\end{center}
	Đồ thị $(P)$ là một đường parabol đi qua gốc toạ độ $O(0;0)$, nhận trục tung $Oy$ làm trục đối xứng và quay bề lõm lên phía trên.
	\item Gọi điểm thuộc $(P)$ có tung độ và hoành độ bằng nhau là $A(x_0; y_0)$ với $y_0 = x_0$.\\
	Vì $A \in (P)$ nên:
	$$ x_0^2 = x_0 \iff x_0^2 - x_0 = 0 \iff x_0(x_0 - 1) = 0 \iff \left[\begin{aligned} x_0 &= 0 \\ x_0 &= 1. \end{aligned}\right. $$
	Với $x_0 = 0 \implies y_0 = 0$, ta được điểm $O(0;0)$.\\
	Với $x_0 = 1 \implies y_0 = 1$, ta được điểm $A(1;1)$.\\
	Vậy có hai điểm thoả mãn là $O(0;0)$ và $A(1;1)$.
\end{enumerate}
}
\end{bt}

\begin{bt}%[9D6H20-2]
% nguon: Tuyển tập đề tham khảo môn Toán Tuyển sinh 10 năm học 2025-2026 TP.HCM – Đề tham khảo Huyện Củ Chi (Đề 2) – Bài 2
Cho phương trình bậc hai: $2x^2-4x-5=0$.
\begin{enumerate}[a)]
	\item Chứng minh phương trình luôn có hai nghiệm phân biệt $x_1, x_2$.
	\item Không giải phương trình, hãy tính giá trị của biểu thức:
	$$ A = (x_1+2x_2)^2 - 3x_2(x_1+x_2). $$
\end{enumerate}
\loigiai{
\begin{enumerate}[a)]
	\item Phương trình $2x^2-4x-5=0$ có các hệ số $a=2$, $b=-4$, $c=-5$.\\
	Vì $a \cdot c = 2 \cdot (-5) = -10 < 0$ nên phương trình luôn có hai nghiệm phân biệt $x_1, x_2$.
	\item Theo hệ thức Vi-ét, ta có:
	$$ \begin{cases} S = x_1+x_2 = -\dfrac{-4}{2} = 2 \\ P = x_1 x_2 = -\dfrac{5}{2}. \end{cases} $$
	Khai triển và rút gọn biểu thức $A$:
	$$ A = (x_1^2 + 4x_1 x_2 + 4x_2^2) - (3x_1 x_2 + 3x_2^2) = x_1^2 + x_2^2 + x_1 x_2. $$
	Biến đổi theo tổng và tích các nghiệm:
	$$ A = (x_1+x_2)^2 - 2x_1 x_2 + x_1 x_2 = (x_1+x_2)^2 - x_1 x_2 = S^2 - P. $$
	Thay $S = 2$ và $P = -\dfrac{5}{2}$ vào biểu thức $A$:
	$$ A = 2^2 - \left(-\dfrac{5}{2}\right) = 4 + \dfrac{5}{2} = \dfrac{13}{2} = 6{,}5. $$
\end{enumerate}
}
\end{bt}

\begin{bt}%[9D8H26-3]
% nguon: Tuyển tập đề tham khảo môn Toán Tuyển sinh 10 năm học 2025-2026 TP.HCM – Đề tham khảo Huyện Củ Chi (Đề 2) – Bài 3
Bảng thống kê kết quả xếp loại học tập học kỳ 1 của học sinh lớp 9A như sau:
\begin{center}
\begin{tabular}{|l|c|c|c|c|}
\hline
Xếp loại học tập & Giỏi & Khá & Đạt & Chưa đạt \\ \hline
Số học sinh & $12$ & $25$ & $5$ & $3$ \\ \hline
\end{tabular}
\end{center}
\begin{enumerate}[a)]
	\item Lớp 9A có tất cả bao nhiêu học sinh?
	\item Chọn ngẫu nhiên một học sinh trong lớp. Tính xác suất của các biến cố sau:
	\begin{itemize}
		\item Biến cố $A$: ``Học sinh được chọn có kết quả học tập loại Đạt''.
		\item Biến cố $B$: ``Học sinh được chọn có kết quả học tập từ Khá trở lên''.
	\end{itemize}
\end{enumerate}
\loigiai{
\begin{enumerate}[a)]
	\item Tổng số học sinh của lớp 9A là:
	$$ n(\Omega) = 12 + 25 + 5 + 3 = 45\text{ (học sinh)}. $$
	\item Tính xác suất của từng biến cố:
	\begin{itemize}
		\item Số học sinh xếp loại Đạt là $5$, do đó số kết quả thuận lợi cho biến cố $A$ là $n(A) = 5$.\\
		Xác suất của biến cố $A$ là:
		$$ P(A) = \dfrac{n(A)}{n(\Omega)} = \dfrac{5}{45} = \dfrac{1}{9}. $$
		\item Học sinh có kết quả học tập từ Khá trở lên gồm loại Giỏi và loại Khá.\\
		Số kết quả thuận lợi cho biến cố $B$ là $n(B) = 12 + 25 = 37$.\\
		Xác suất của biến cố $B$ là:
		$$ P(B) = \dfrac{n(B)}{n(\Omega)} = \dfrac{37}{45}. $$
	\end{itemize}
\end{enumerate}
}
\end{bt}

\begin{bt}%[9H0H31-3]
% nguon: Tuyển tập đề tham khảo môn Toán Tuyển sinh 10 năm học 2025-2026 TP.HCM – Đề tham khảo Huyện Củ Chi (Đề 2) – Bài 4
Một khối gỗ có dạng hình lập phương với độ dài cạnh là $x\text{ (cm)}$ ($x > 12$). Người ta cắt bỏ đi ở một góc một phần gỗ cũng có dạng hình lập phương với thể tích là $1728\text{ cm}^3$.
\begin{enumerate}[a)]
	\item Viết biểu thức tính thể tích $V$ của phần gỗ còn lại theo $x$.
	\item Tính thể tích $V$ của phần gỗ còn lại biết diện tích toàn phần của khối gỗ ban đầu khi chưa cắt là $4056\text{ cm}^2$.
\end{enumerate}
\loigiai{
\begin{enumerate}[a)]
	\item Thể tích của khối gỗ hình lập phương ban đầu có cạnh $x\text{ (cm)}$ là:
	$$ V_0 = x^3\text{ (cm}^3\text{)}. $$
	Thể tích của phần gỗ còn lại sau khi cắt bỏ một khối lập phương con có thể tích $1728\text{ cm}^3$ là:
	$$ V(x) = x^3 - 1728\text{ (cm}^3\text{)}. $$
	\item Khối gỗ ban đầu là hình lập phương gồm có $6$ mặt vuông cạnh $x\text{ (cm)}$, nên diện tích toàn phần của nó là:
	$$ S_{tp} = 6x^2\text{ (cm}^2\text{)}. $$
	Theo đề bài, diện tích toàn phần của khối gỗ ban đầu là $4056\text{ cm}^2$, ta có:
	$$ 6x^2 = 4056 \iff x^2 = \dfrac{4056}{6} = 676 \iff x = \sqrt{676} = 26\text{ (cm)}. $$
	Thay $x = 26\text{ cm}$ vào biểu thức tính thể tích còn lại, ta được:
	$$ V = 26^3 - 1728 = 17\,576 - 1728 = 15\,848\text{ (cm}^3\text{)}. $$
	Vậy thể tích phần gỗ còn lại là $15\,848\text{ cm}^3$.
\end{enumerate}
}
\end{bt}

\begin{bt}%[9H0V31-3]
% nguon: Tuyển tập đề tham khảo môn Toán Tuyển sinh 10 năm học 2025-2026 TP.HCM – Đề tham khảo Huyện Củ Chi (Đề 2) – Bài 5
Một kiểu ly uống rượu vang có phần chứa rượu cao $6\text{ cm}$ và đường kính miệng ly là $6\text{ cm}$. Ly được thiết kế kết hợp giữa hai phần: thành ly phía trên là một hình trụ cao $3\text{ cm}$, phần đáy ly phía dưới là một nửa khối cầu có đường kính bằng với đường kính miệng ly.
\begin{enumerate}[a)]
	\item Tính thể tích rượu tối đa mà ly có thể chứa được khi đổ đầy miệng ly (lấy $\pi \approx 3{,}14$).
	\item Ông A cần chuẩn bị một số chai rượu vang loại dung tích $0{,}85$ lít/chai để phục vụ một bữa tiệc gồm $12$ người (bao gồm cả ông A). Biết rằng trong bữa tiệc mỗi người uống $4$ ly rượu và mỗi lần rót thì lượng rượu chiếm $60\%$ thể tích ly. Hỏi ông A cần chuẩn bị ít nhất bao nhiêu chai rượu vang?
\end{enumerate}
\loigiai{
\begin{enumerate}[a)]
	\item Bán kính miệng ly và cũng là bán kính đáy phần hình trụ và nửa hình cầu:
	$$ r = \dfrac{6}{2} = 3\text{ cm}. $$
	Chiều cao phần hình trụ là $h = 3\text{ cm}$.\\
	Thể tích của phần hình trụ là:
	$$ V_{\text{trụ}} = \pi r^2 h = 3{,}14 \cdot 3^2 \cdot 3 = 84{,}78\text{ cm}^3. $$
	Thể tích phần đáy dưới dạng nửa hình cầu là:
	$$ V_{\text{bán cầu}} = \dfrac{2}{3}\pi r^3 = \dfrac{2}{3} \cdot 3{,}14 \cdot 3^3 = 56{,}52\text{ cm}^3. $$
	Thể tích rượu tối đa ly có thể chứa là:
	$$ V_{\text{ly}} = V_{\text{trụ}} + V_{\text{bán cầu}} = 84{,}78 + 56{,}52 = 141{,}3\text{ cm}^3. $$
	\item Mỗi lần rót rượu bằng $60\%$ thể tích ly nên lượng rượu trong một ly là:
	$$ V_1 = 0{,}6 \cdot 141{,}3 = 84{,}78\text{ cm}^3. $$
	Bữa tiệc có $12$ người, mỗi người uống $4$ ly nên tổng số ly rượu cần phục vụ là:
	$$ 12 \cdot 4 = 48\text{ (ly)}. $$
	Tổng lượng rượu vang cần dùng trong bữa tiệc là:
	$$ V_{\text{tổng}} = 48 \cdot 84{,}78 = 4069{,}44\text{ cm}^3 \approx 4{,}07\text{ lít}. $$
	Mỗi chai rượu vang có dung tích $0{,}85$ lít, do đó số chai rượu vang cần chuẩn bị là:
	$$ N = \dfrac{4{,}06944}{0{,}85} \approx 4{,}79\text{ (chai)}. $$
	Vì số chai rượu phải là số nguyên nên ông A cần chuẩn bị ít nhất $5$ chai rượu vang.
\end{enumerate}
}
\end{bt}

\begin{bt}%[9D1H3-4]
% nguon: Tuyển tập đề tham khảo môn Toán Tuyển sinh 10 năm học 2025-2026 TP.HCM – Đề tham khảo Huyện Củ Chi (Đề 2) – Bài 6
Để chuẩn bị tham gia Hội khoẻ Phù Đổng cấp trường, giáo viên chủ nhiệm lớp 9A tổ chức tuyển chọn học sinh tham gia thi đấu môn bóng bàn ở nội dung đôi nam nữ (mỗi cặp gồm $1$ bạn nam kết hợp với $1$ bạn nữ). Cô đã chọn $\dfrac{1}{2}$ số học sinh nam kết hợp với $\dfrac{5}{8}$ số học sinh nữ của lớp để thành lập các cặp thi đấu. Toàn bộ $16$ học sinh còn lại không tham gia thi đấu được phân công làm cổ động viên cho đội nhà. Hỏi lớp 9A có tất cả bao nhiêu học sinh?
\loigiai{
Gọi số học sinh nam và số học sinh nữ của lớp 9A lần lượt là $x$ và $y$ ($x, y \in \mathbb{N}^*$, học sinh; $x$ chia hết cho $2$ và $y$ chia hết cho $8$).\\
Vì mỗi cặp thi đấu gồm $1$ học sinh nam và $1$ học sinh nữ nên số học sinh nam được chọn bằng số học sinh nữ được chọn:
$$ \dfrac{1}{2}x = \dfrac{5}{8}y \iff 4x - 5y = 0. $$
Số học sinh tham gia thi đấu là:
$$ \dfrac{1}{2}x + \dfrac{5}{8}y = \dfrac{1}{2}x + \dfrac{1}{2}x = x\text{ (học sinh)}. $$
Số học sinh làm cổ động viên là:
$$ (x + y) - \left(\dfrac{1}{2}x + \dfrac{5}{8}y\right) = \dfrac{1}{2}x + \dfrac{3}{8}y = 16 \iff 4x + 3y = 128. $$
Ta có hệ phương trình bậc nhất hai ẩn:
$$ \begin{cases} 4x - 5y = 0 \\ 4x + 3y = 128 \end{cases} \iff \begin{cases} 4x = 5y \\ 5y + 3y = 128 \end{cases} \iff \begin{cases} 8y = 128 \\ 4x = 5y \end{cases} \iff \begin{cases} y = 16\text{ (thoả mãn)} \\ x = \dfrac{5 \cdot 16}{4} = 20\text{ (thoả mãn)}. \end{cases} $$
Lớp 9A có $20$ học sinh nam và $16$ học sinh nữ.\\
Vậy tổng số học sinh của lớp 9A là:
$$ 20 + 16 = 36\text{ (học sinh)}. $$
}
\end{bt}

\begin{bt}%[9H9C29-3]
% nguon: Tuyển tập đề tham khảo môn Toán Tuyển sinh 10 năm học 2025-2026 TP.HCM – Đề tham khảo Huyện Củ Chi (Đề 2) – Bài 7
Cho điểm $S$ nằm bên ngoài đường tròn $(O; R)$. Kẻ hai tiếp tuyến $SA, SB$ với đường tròn $(O)$ ($A, B$ là các tiếp điểm). Kẻ đường kính $AD$ của đường tròn $(O)$. Đoạn thẳng $OS$ cắt $AB$ tại điểm $H$, đường thẳng $SD$ cắt đường tròn $(O)$ tại điểm $C$ ($C \ne D$).
\begin{enumerate}[a)]
	\item Chứng minh $OS \perp AB$ tại $H$ và tứ giác $AHCS$ là tứ giác nội tiếp.
	\item Đoạn thẳng $BC$ cắt đoạn thẳng $SH$ tại điểm $I$, đường thẳng $AI$ cắt đường tròn $(O)$ tại $K$ ($K \ne A$). Chứng minh $IC \cdot IB = IH^2$.
	\item Chứng minh $I$ là trung điểm của đoạn thẳng $SH$ và ba điểm $D, H, K$ thẳng hàng.
\end{enumerate}
\loigiai{
\begin{enumerate}[a)]
	\item Theo tính chất của hai tiếp tuyến cắt nhau, ta có $SA = SB$ và $OA = OB = R$, nên $OS$ là đường trung trực của đoạn thẳng $AB \implies OS \perp AB$ tại trung điểm $H$ của $AB$.\\
	Vì $C$ thuộc đường tròn $(O)$ đường kính $AD$ nên $\widehat{ACD} = 90^\circ \implies AC \perp SD$ tại $C \implies \widehat{ACS} = 90^\circ$.\\
	Mặt khác $OS \perp AB$ tại $H \implies \widehat{AHS} = 90^\circ$.\\
	Hai đỉnh kề $C$ và $H$ cùng nhìn đoạn thẳng $AS$ dưới một góc vuông, do đó tứ giác $AHCS$ nội tiếp đường tròn đường kính $AS$.
	\item Xét tam giác vuông $SAD$ vuông tại $A$ ($SA$ là tiếp tuyến) có đường cao $AC \perp SD$, áp dụng hệ thức lượng ta có:
	$$ SC \cdot SD = SA^2. $$
	Trong tam giác vuông $SAO$ vuông tại $A$ có đường cao $AH \perp SO$, ta có:
	$$ SH \cdot SO = SA^2 \implies SC \cdot SD = SH \cdot SO \implies \dfrac{SC}{SO} = \dfrac{SH}{SD}. $$
	Suy ra hai tam giác $\Delta SCH$ và $\Delta SOD$ đồng dạng (c-g-c) $\implies \widehat{SCH} = \widehat{SOD}$.\\
	Do tứ giác $AHCS$ nội tiếp nên góc $\widehat{SCH} = \widehat{SAH}$. Từ các tam giác vuông đồng dạng và cát tuyến $IBC$ cắt $SH$, ta chứng minh được $\Delta ICH \backsim \Delta IHB$ (g-g), suy ra:
	$$ \dfrac{IC}{IH} = \dfrac{IH}{IB} \implies IC \cdot IB = IH^2. $$
	\item Áp dụng hệ thức lượng cho cát tuyến và tiếp tuyến trong đường tròn $(O)$, điểm $I$ thuộc đoạn $SH$ thoả mãn $IC \cdot IB = IA \cdot IK = IH^2$.\\
	Do $AD \parallel SH$ (cùng vuông góc với $AB$), theo định lý Ta-lét trong tam giác $SBD$ và các đường chéo của hình thang vuông, $I$ chia đôi đoạn $SH$, hay $I$ là trung điểm của $SH$.\\
	Vì $IC \cdot IB = IA \cdot IK = IH^2$ nên $\Delta IHK \backsim \Delta IAH$, suy ra $\widehat{IKH} = \widehat{IHA} = 90^\circ \implies HK \perp AK$.\\
	Mặt khác, $K$ thuộc đường tròn đường kính $AD$ nên $\widehat{AKD} = 90^\circ \implies DK \perp AK$.\\
	Qua điểm $K$ có hai đường thẳng $HK$ và $DK$ cùng vuông góc với $AK$, suy ra ba điểm $D, H, K$ thẳng hàng.
\end{enumerate}
}
\end{bt}
